\chapter{字体结构}\label{fonts}

\section{字体表}

\topicindex {fonts}
\topicindex {fonts+tables}

\TEX\ 的所有字体在 \LUA\ 代码中都表示为表（table），在内部则表示为
\CCODE~结构体。下表中的所有键，如果出现在 \cbk {define_font} 回调所返回
的表中，就会被保存到内部字体结构体中；如果没有定义 \cbk {define_font}
回调，则由常规的 \TFM|/|\VF\ 读取例程产生的键也会被保存。

列 ``\VF 表示该键由 \type {font.read_vf()} 例程创建，列
``\TFM 表示该键由 \type {font.read_tfm()} 例程创建，列 ``used
表示 \LUATEX\ 引擎本身是否会使用该键。表中的顶层键如下：

\begin{longtable}{l c c c l p{5cm}}
\toprule 键                     & vf  & tfm & used & 值类型   & 说明 \\
\midrule
\type{name}             & yes & yes & yes  & string     & 度量（文件）名 \\
\type{area}             & no  & yes & yes  & string     & （目录）位置，通常为空 \\
\type{used}             & no  & yes & yes  & boolean    & 指示是否被使用（初始值：false） \\
\type{characters}       & yes & yes & yes  & table      & 此字体中已定义的字形 \\
\type{checksum}         & yes & yes & no   & number     & 默认值：0 \\
\type{designsize}       & no  & yes & yes  & number     & 期望尺寸（默认值：655360 == 10pt） \\
\type{direction}        & no  & yes & yes  & number     & 默认值：0 \\
\type{encodingbytes}    & no  & no  & yes  & number     & 默认值：取决于 \type {format} \\
\type{encodingname}     & no  & no  & yes  & string     & 编码名称 \\
\type{fonts}            & yes & no  & yes  & table      & 本地使用的字体 \\
\quad\type{psname}           & no  & no  & yes  & string     & 这是输入字体源中的 \POSTSCRIPT\ 字体名，
                                                                        用作 \PDF\ 输出中的字体名标识符。它必须是合法
                                                                        字符串（例如不能包含空格之类的字符），因为后
                                                                        端不会做清理。这样就把完全的控制权交给了加载器。 \\
\quad\type{fullname}         & no  & no  & yes  & string     & 输出字体名，当未设置 \type {psname} 时在
                                                                        \PDF\ 输出中用作后备 \\
\quad\type{subfont}          & no  & no  & yes  & number     & 默认值：0，多字体（\type {ttc}）文件中的
                                                                        子字体索引 \\
\type{header}           & yes & no  & no   & string     & 头部注释（如果有） \\
\type{hyphenchar}       & no  & no  & yes  & number     & 默认值：\TEX\ 的 \prm {hyphenchar} \\
\type{parameters}       & no  & yes & yes  & hash       & 默认值：7 个参数，全部为零 \\
\type{size}             & no  & yes & yes  & number     & 所需的缩放（默认与 designsize 相同） \\
\type{skewchar}         & no  & no  & yes  & number     & 默认值：\TEX\ 的 \prm {skewchar} \\
\type{type}             & yes & no  & yes  & string     & 此字体的基本类型 \\
\type{format}           & no  & no  & yes  & string     & 磁盘格式类型 \\
\type{embedding}        & no  & no  & yes  & string     & \PDF\ 嵌入方式  \\
\type{filename}         & no  & no  & yes  & string     & 磁盘上的字体文件名 \\
\quad\type{tounicode}        & no  & yes & yes  & number     & 当此项设为~1 时，\LUATEX\ 假定字体中
                                                                        存在逐字形的 tounicode 条目。 \\
\quad\type{stretch}          & no  & no  & yes  & number     & 来自 \lpr {expandglyphsinfont} 的
                                                                        ``stretch 值 \\
\quad\type{shrink}           & no  & no  & yes  & number     & 来自 \lpr {expandglyphsinfont} 的
                                                                        ``shrink 值 \\
\quad\type{step}             & no  & no  & yes  & number     & 来自 \lpr {expandglyphsinfont} 的
                                                                        ``step 值 \\
\type{expansion_factor} & no  & no  & no   & number     & 已扩展字体的实际扩展系数 \\
\type{attributes}       & no  & no  & yes  & string     & \orm {pdffontattr} \\
\quad\type{cache}            & no  & no  & yes  & string     & 此键控制 \TEX\ 端对 \LUA\ 表的缓存：
                                                                        \type {yes} 表示使用传递给 \LUATEX\ 的表的引用
                                                                        （这是默认值）；\type {no} 表示不保存表引用，
                                                                        也不为该字体缓存任何 \LUA\ 数据；而
                                                                        \type {renew} 表示不保存表引用，但保存对
                                                                        首次访问字体某个字段时所创建表的引用。 \\
\quad\type{nomath}           & no  & no  & yes  & boolean    & 此键可为文本字体带来小幅提速。如果存在且
                                                                        为 true，\LUATEX\ 将不检查字符条目中的数学
                                                                        专用键。 \\
\quad\type{oldmath}          & no  & no  & yes  & boolean    & 此键把字体标记为旧式 \TEX\ 数学字体，
                                                                        并禁用 \OPENTYPE\ 代码路径。 \\
\quad\type{slant}            & no  & no  & yes  & number     & 此参数使字体倾斜，作用与 \TYPEONE\ 字体
                                                                        映射文件中的 \type {SlantFont} 相同。 \\
\quad\type{extend}           & no  & no  & yes  & number     & 此参数在水平方向上缩放字体，作用与
                                                                        \TYPEONE\ 字体映射文件中的 \type {ExtendFont}
                                                                        相同。 \\
\quad\type{squeeze}          & no  & no  & yes  & number     & 此参数在垂直方向上缩放字体，映射文件中
                                                                        没有对应项。 \\
\quad\type{width}            & no  & no  & yes  & number     & 后端会注入设置笔宽的 \PDF\ 运算符。该值
                                                                        （与 \TEX\ 中惯例一样）除以 1000。它与
                                                                        \type {mode} 配合使用。 \\
\quad\type{mode}             & no  & no  & yes  & number     & 后端会注入与绘制模式相关的 \PDF\ 运算符：
                                                                        0~为填充，1~为轮廓，2~既绘制又填充，
                                                                        3~完全不绘制。 \\
\bottomrule
\end{longtable}

\type {cache} 选项中保存的引用是线程局部的，因此使用协程时要小心：如果
表已在一个线程中缓存，而你从另一个线程引用它，就会抛出错误。

键 \type {name} 总是必需的。键 \type {stretch}、\type {shrink}、\type
{step} 只有一起使用时才有意义：它们可用来替代加载后的
\lpr {expandglyphsinfont} 命令。\LUATEX\ 不支持 \type {auto_expand}
选项。事实上，那些创建扩展或突出副本的原语大概只在与传统字体配合时有用，
因为所有这些额外的 \OPENTYPE\ 属性都被排除在外。\type {expansion_factor}
是一个可以出现在 \type {font.fonts} 中某个字体内部的值。它是由字体扩展
算法自动生成的字体的实际扩展系数（取值在 \type {-shrink} 与
\type {stretch} 之间，步长为 \type {step}）。

\type {subfont} 参数可用于指定 \type {ttc} 字体中的子字体。给出该参数时，
它会取代 \type {psname} 与 \type {fullname} 的组合。第一个子字体的编号
为~1。值为零表示使用名字进行查找。

由于我们把扩展的实际状态随每个字形一起存储，并且没有专门的字体实例，
因此可以在构造行之前更改一些与字体相关的参数，例如：

\begin{typing}
font.setexpansion(font.current(),100,100,20)
\end{typing}

这主要是为了实验（或用 \LUA\ 编写的优化例程），因此没有对应的原语。

键 \type {attributes} 可用于在 \PDF\ 文件中设置字体属性。键 \type {used}
由引擎在字体被实际使用时置位，这会确保字体的定义被写入输出文件（\DVI\
或 \PDF）。\TFM\ 读取器把它设为 false。\type {direction} 是一个数字，
表示此字体的 ``正常 方向。共有十六种可能：

\begin{longtable}{c c c c c c c c}
\toprule \# & 方向 & \# & 方向 & \# & 方向 & \# & 方向 \\
\midrule
0 & LT  &  4 & RT  &  8 & TT  & 12 & BT  \\
1 & LL  &  5 & RL  &  9 & TL  & 13 & BL  \\
2 & LB  &  6 & RB  & 10 & TB  & 14 & BB  \\
3 & LR  &  7 & RR  & 11 & TR  & 15 & BR  \\
\bottomrule
\end{longtable}

这些是 \OMEGA\ 风格的方向缩写：第一个字符表示字符字形的 ``首 边
（在书写方向上最先看到的边），第二个字符表示 ``顶 边。请记住，
\LUATEX\ 的方向模型略有不同，因此这些值实际上并不派什么用场。

\type {parameters} 是一个键类型混合的哈希。它有七个可能的字符串键，
以及若干个整数索引（从 8 开始向上编号）。这七个字符串实际上用来替代
最下面的七个索引，因为这样的用户界面更友好。

这些名字及其内部重映射如下：

\begin{longtable}{l c}
\toprule 名称                  & 重映射 \\
\midrule
\type {slant}         & 1 \\
\type {space}         & 2 \\
\type {space_stretch} & 3 \\
\type {space_shrink}  & 4 \\
\type {x_height}      & 5 \\
\type {quad}          & 6 \\
\type {extra_space}   & 7 \\
\bottomrule
\end{longtable}

键 \type {type}、\type {format}、\type {embedding}、\type {fullname} 和
\type {filename} 用于把 \OPENTYPE\ 字体嵌入结果 \PDF\ 中。

\type {characters} 表是一个字符哈希列表，以整数作为索引。该整数就是
\TEX\ 用来识别此字符的 ``内部代码。

还可以使用两个非常特殊的字符串索引：\type {left_boundary} 是一个虚拟
字符，其连字和字距用于处理单词边界。\type {right_boundary} 与之类似，
但（目前）实际上并没有派什么用场。

每个字符哈希本身也是一个哈希。例如，下面是字体 \type {cmr10 at 10pt}
中字符 ``f（十进制 102）的例子。表示尺寸的数字单位是缩放点（sp）。

\begin{typing}
[102] = {
    ["width"]  = 200250,
    ["height"] = 455111,
    ["depth"]  = 0,
    ["italic"] = 50973,
    ["kerns"]  = {
        [63] = 50973,
        [93] = 50973,
        [39] = 50973,
        [33] = 50973,
        [41] = 50973
    },
    ["ligatures"] = {
        [102] = { ["char"] = 11, ["type"] = 0 },
        [108] = { ["char"] = 13, ["type"] = 0 },
        [105] = { ["char"] = 12, ["type"] = 0 }
    }
}
\end{typing}

字符哈希中可以出现以下顶层键：

\begin{longtable}{l c c c l p{5cm}}
\toprule 键                     & vf  & tfm & used  & 类型    & 说明 \\
\midrule
\type{width}            & yes & yes & yes   & number  & 字符宽度，单位 sp（默认 0） \\
\type{height}           & no  & yes & yes   & number  & 字符高度，单位 sp（默认 0） \\
\type{depth}            & no  & yes & yes   & number  & 字符深度，单位 sp（默认 0） \\
\type{italic}           & no  & yes & yes   & number  & 字符的斜体校正，单位 sp（默认零） \\
\type{top_accent}       & no  & no  & maybe & number  & 字符的上重音对齐位置，单位 sp（默认零） \\
\type{bot_accent}       & no  & no  & maybe & number  & 字符的下重音对齐位置，单位 sp（默认零） \\
\type{left_protruding}  & no  & no  & maybe & number  & 字符的 \lpr {lpcode} \\
\type{right_protruding} & no  & no  & maybe & number  & 字符的 \lpr {rpcode} \\
\type{expansion_factor} & no  & no  & maybe & number  & 字符的 \lpr {efcode} \\
\type{tounicode}        & no  & no  & maybe & string  & 字符对应的 \UNICODE\ 字符，采用 \UTF|-|16BE 十六进制格式 \\
\type{unicode}           & yes & no  & maybe & table   & 字符对应的 \UNICODE\ 字符，以 \UNICODE\ 码点（整数）数组形式给出 \\
\type{next}             & no  & yes & yes   & number  & ``下一个更大 的字符索引 \\
\type{extensible}       & no  & yes & yes   & table   & 可伸缩配方的组成部分 \\
\type{vert_variants}    & no  & no  & yes   & table   & 竖排变体集合的组成部分 \\
\type{horiz_variants}   & no  & no  & yes   & table   & 横排变体集合的组成部分 \\
\type{kerns}            & no  & yes & yes   & table   & 字距调整信息 \\
\type{ligatures}        & no  & yes & yes   & table   & 连字信息 \\
\type{commands}         & yes & no  & yes   & array   & 虚拟字体命令 \\
\type{name}             & no  & no  & no    & string  & 字符的（\POSTSCRIPT）名字 \\
\type{index}            & no  & no  & yes   & number  & （\OPENTYPE\ 或 \TRUETYPE）字体的字形索引 \\
\type{used}             & no  & yes & yes   & boolean & 已被排版过（默认：false） \\
\type{mathkern}         & no  & no  & yes   & table   & 数学切入（cut-in）规格 \\
\bottomrule
\end{longtable}

\type {top_accent}、\type {bot_accent} 和 \type {mathkern} 的值仅用于
数学重音和上标的放置，详见第~\ref{math}~节。\type {left_protruding}
和 \type {right_protruding} 的值仅在 \lpr {protrudechars} 非零时使用。
是否使用 \type {expansion_factor} 取决于字体的全局扩展设置以及
\lpr {adjustspacing} 的值。

\type {tounicode} 的用法如下：如果字体在顶层指定了
\type {tounicode=1}，那么 \LUATEX\ 会根据字符级别的 \type {tounicode}
字符串（在有提供的地方）为 \PDF\ 字体（或字体子集）构造
\type {/ToUnicode} 条目。如果某个字符没有合理的 \UNICODE\ 对应字符，
就不要提供字符串（不要给空字符串）。

如果没有设置字体级别的 \type {tounicode}，\LUATEX\ 将根据你使用的
\TEX\ 码点来构建 \type {/ToUnicode}，任何字符级别的 \type {tounicodes}
都会被忽略。字符串格式正是 Adobe \CMAP\ 文件所期望的格式（十六进制编码
的 \UTF-16BE），只是去掉了外层的尖括号。例如，\type {fi} 连字的
\type {tounicode} 应为 \type {00660069}。如果你传入的是数字，转换会
自动完成。

数学字符可以有一个 \type {next} 字段，指向下一个更大的字形。不过，
如果同时存在 \type {extensible}，它会覆盖 \type {next}。而
\type {extensible} 字段又会被 \OPENTYPE\ 版本的 \type {vert_variants}
覆盖。\type {extensible} 表非常简单：

\begin{longtable}{l l p{5cm}}
\toprule 键        & 类型   & 说明                \\
\midrule
\type{top} & number & 顶部字符索引        \\
\type{mid} & number & 中部字符索引        \\
\type{bot} & number & 底部字符索引        \\
\type{rep} & number & 可重复的字符索引    \\
\bottomrule
\end{longtable}

\type {horiz_variants} 和 \type {vert_variants} 是组件数组。每个组件
本身又是一个最多含五个键的哈希：

\begin{longtable}{l l p{5cm}}
\toprule 键             & 类型   & 说明 \\
\midrule
\type{glyph}    & number & 字符索引。注意这是编码数字，不是名字。 \\
\type{extender} & number & 若此部分可重复则为一（1），否则为零（0）。 \\
\type{start}    & number & 起始侧的最大重叠量（单位：缩放点）。 \\
\type{end}      & number & 结束侧的最大重叠量（单位：缩放点）。 \\
\quad\type{advance}  & number & 此项的总前进宽度。它可以为零或缺失，
                                   此时使用字符 \type {component} 对应字形的
                                   自然尺寸。 \\
\bottomrule
\end{longtable}

\type {kerns} 表是一个以字符索引为键的哈希（``字符索引 定义为
非负整数或字符串值 \type {right_boundary}），其值为要应用的字距调整量，
单位为缩放点。

\type {ligatures} 表是一个以字符索引为键的哈希（``字符索引 定义为
非负整数或字符串值 \type {right_boundary}），其值又是一个小哈希，
含两个字段：

\begin{longtable}{l l p{5cm}}
\toprule 键         & 类型   & 说明 \\
\midrule
\type{type} & number & 此连字命令的类型，默认 0 \\
\type{char} & number & 结果连字的字符索引 \\
\bottomrule
\end{longtable}

连字中的 \type {char} 字段是必需的。连字内部的 \type {type} 字段是
\TEX\ 支持的八种连字类型之一的数字或字符串值。当 \TEX\ 插入一个新连字时，
它把新字形放在左右两个字形的中间。原来的左右字形可以选择保留；当至少
保留其中之一时，还可以把新的 ``插入点 向前移动一位或两位。最终
位于插入点右侧的字形将成为下一个 ``左 字形。

\begin{longtable}{l c l l}
\toprule 文字形式（Knuth）      & 数字   & 字符串        & 结果      \\
\midrule
\type{l + r =: n}     & 0      & \type{=:}     & \type{|n}   \\
\type{l + r =:| n}    & 1      & \type{=:|}    & \type{|nr}  \\
\type{l + r |=: n}    & 2      & \type{|=:}    & \type{|ln}  \\
\type{l + r |=:| n}   & 3      & \type{|=:|}   & \type{|lnr} \\
\type{l + r  =:|> n}  & 5      & \type{=:|>}   & \type{n|r}  \\
\type{l + r |=:> n}   & 6      & \type{|=:>}   & \type{l|n}  \\
\type{l + r |=:|> n}  & 7      & \type{|=:|>}  & \type{l|nr} \\
\type{l + r |=:|>> n} & 11     & \type{|=:|>>} & \type{ln|r} \\
\bottomrule
\end{longtable}

默认值为~0，可以省略。它表示一个 ``常规 连字，即连字替换原来的
两个字形。在此表中，\type {|} 表示最终的插入点。

\type {commands} 数组在下面说明。

\section{真实字体}

\topicindex {fonts+real}
\topicindex {fonts+virtual}

一个 \TEX\ 字体是否是应该写入 \PDF\ 文档的 ``真实 字体，由顶层字体
结构中的 \type {type} 值决定。如果该值为 \type {real}，则它是一个真正的
字体，嵌入机制会尝试把所需的字体对象定义添加到 \PDF\ 中。\type {type}
的取值如下：

\begin{longtable}{l p{5cm}}
\toprule 取值          & 说明            \\
\midrule
\type{real}    & 这是一个基础字体    \\
\type{virtual} & 这是一个虚拟字体 \\
\bottomrule
\end{longtable}

要采取的操作取决于若干不同的变量：

\begin{itemize}
\item 所用字体是否适合 8 位编码方案。对于通过 \TFM\ 文件通信的传统
    \TEX\ 字体，答案是肯定的。

\item 磁盘字体文件的类型，例如位图文件，或者轮廓型 \TYPEONE、\TRUETYPE\
    或 \OPENTYPE\ 字体。

\item 请求的嵌入级别，尽管大多数情况下嵌入的是字符子集。完全不嵌入的
    时代（至少在我们看来）基本上已经过去了。

\end{itemize}

使用 8 位编码向量之外任何方案的字体，必须以不同方式写入 \PDF。当字体表
的 \type {encodingbytes} 设为~2 时，它就是宽字体，其他情况都不是。
对于通过 \LUA\ 加载的 \OPENTYPE\ 和 \TRUETYPE\ 字体，值~2 是默认值。
对于 \TYPEONE\ 字体，必须显式把 \type {encodingbytes} 设为~2。对于
\PK\ 位图字体，完全不支持宽字体编码。

如果不需要特殊处理，\LUATEX\ 会回退到 \PDFTEX\ 和 \DVIPS\ 所使用的
基于映射文件的方案，从而透明地支持旧版字体。如果使用了 ``宽
字体，新的子系统就会启动，字体结构中必须存在一些额外字段。在这种情况下，
\LUATEX\ 完全不使用映射文件。这些额外字段是：\type {format}、
\type {embedding}、\type {fullname}、\type {cidinfo}（如前所述）、
\type {filename}，以及各个字符中的 \type {index} 键。

\type {format} 变量可取以下值。\type {type3} 字体仅为向后兼容而提供，
不支持新的宽编码选项。

\begin{longtable}{l p{5cm}}
\toprule 取值           & 说明                                                      \\
\midrule
\type{type1}    & 这是一个 \POSTSCRIPT\ \TYPEONE\ 字体                       \\
\type{type3}    & 这是一个位图（\PK）字体                                   \\
\type{truetype} & 这是一个 \TRUETYPE\ 或基于 \TRUETYPE\ 的 \OPENTYPE\ 字体 \\
\type{opentype} & 这是一个基于 \POSTSCRIPT\ 的 \OPENTYPE\ 字体                \\
\bottomrule
\end{longtable}

\type {embedding} 变量的合法取值如下：

\begin{longtable}{l p{5cm}}
\toprule 取值         & 说明                           \\
\midrule
\type{no}     & 完全不嵌入该字体                 \\
\type{subset} & 包含字体并尝试子集化             \\
\type{full}   & 完整包含此字体                   \\
\bottomrule
\end{longtable}

其余字段的用法如下。\type {fullname} 将作为 \POSTSCRIPT|/|\PDF\ 字体名。
\type {cidinfo} 将用作字符集，即 CID 的 \type {/Ordering} 和
\type {/Registry} 键。\type {filename} 指向实际的字体文件。如果你在
\type {filename} 中给出完整路径，或者文件就在当前目录下，\LUATEX\ 的
运行效率会稍高一些，因为那种情况下不必重新运行 \type {find_*_file}
回调。

请注意：在一个文档中混用新旧字体时，可能会产生 \POSTSCRIPT\ 名字冲突，
从而导致打印错误。出现这种情况时，必须把字体的 \type {fullname} 改成
更唯一的名字。

排版字符串以宽格式写出，每个字形使用 2~字节，以字符信息中的
\type {index} 键作为值。整体效果就像是采用了基于数字的编码，而不是
传统的（\POSTSCRIPT）基于名字的重编码。获取 \TYPEONE\ 字体正确
\type {index} 编号的一种方法是通过 \type {fontloader.open} 加载字体，
并把表索引用作 \type {index} 字段。

为了确保最终文档的剪切和粘贴正常工作，最好强制提供 \type {tounicode}
向量。并非所有 \PDF\ 阅读器都能正确处理这一点，因此请以 \ACROBAT\
为参照。

\section{虚拟字体}\label{virtualfonts}

\subsection{结构}

\topicindex {fonts+virtual}

如果你想让 \LUATEX\ 把 \cbk {define_font} 返回的表当作虚拟字体处理，
必须采取以下步骤：

\begin{itemize}
\item 把顶层键 \type {type} 设为 \type {virtual}。在大多数情况下这是可选的，
    因为我们反正会查看 \type {commands} 条目。

\item 确保 \type {fonts} 中至少有一个有效条目（见下文），不过在该表缺失时，
    较新版本的 \LUATEX\ 会添加一个默认条目。

\item 在相关字符上添加 \type {commands} 数组。虚拟字符本身也可以指向其他
    虚拟字符，但要小心嵌套，因为这可能造成循环并使栈溢出（而栈溢出通常
    也意味着出错了）。

\end{itemize}

顶层 \type {type} 键取特定值 \type {virtual} 会触发对表中其余虚拟字体
专用字段的处理，但仅仅存在 'type' 键就足以阻止 \LUATEX\ 自行查找虚拟
字体。这也可以 ``反向 运用：如果你完全确定某个字体不是虚拟字体，
把 \type {real} 赋给 \type {type} 会阻止 \LUATEX\ 查找虚拟字体文件，
从而省去一次磁盘搜索。这只在加载 \TFM\ 文件时才有意义。

\type {fonts} 是一个（带索引的）\LUA\ 表。其值本身是含一两个键的哈希，
每个条目表示虚拟字体中的一个基础字体。如果你的字体要引用自身，可以使用
\type {font.nextid()} 函数，它返回下一个将要定义的字体的索引，而那
很可能就是当前正在定义的字体。因此，一个表看起来像这样：

\begin{typing}
fonts = {
  { name = "ptmr8a", size = 655360 },
  { name = "psyr", size = 600000 },
  { id = 38 }
}
\end{typing}

此虚拟字体中第一个被引用的字体（索引~1）是按 10pt 加载的
\type {ptmr8a}，第二个是按略大于 9pt 加载的 \type {psyr}。第三个是先前
定义过的字体，在 \LUATEX\ 中已知为字体 id~38。数组索引编号由属于各字符
的字符命令定义使用。

\type {commands} 数组是一个哈希，其中每一项又是一个小数组，第一项表示
命令，其余各项是该命令的参数。允许的命令及其参数如下：

\begin{longtable}{l l l p{5cm}}
\toprule 命令        & 参数个数 & 类型      & 说明 \\
\midrule
\type{font}    & 1         & number    & 从本地 \type {fonts} 表中选择一个新字体 \\
\type{char}    & 1         & number    & 用当前字体排版此编号字符，并向右移动该字符的宽度 \\
\type{node}    & 1         & node      & 输出此节点（列表），并向右移动该列表的宽度\\
\type{slot}    & 2         & 2 numbers & font 与 char 命令组合的快捷方式\\
\type{push}    & 0         &           & 保存当前位置\\
\type{nop}     & 0         &           & 什么也不做 \\
\type{pop}     & 0         &           & 弹出位置 \\
\type{rule}    & 2         & 2 numbers & 输出一个 $ht*wd$ 的矩形，并向右移动。 \\
\type{down}    & 1         & number    & 在页面上向下移动 \\
\type{right}   & 1         & number    & 在页面上向右移动 \\
\type{special} & 1         & string    & 输出一个 \prm {special} 命令 \\
\type{pdf}     & 2         & 2 strings & 输出一个 \PDF\ 字面量，第一个字符串取 \type {origin}、
                                                   \type {page}、\type {text}、\type {font}、\type {direct}
                                                   或 \type {raw} 之一；如果只有一个字符串，则假定为
                                                   \type {origin} \\
\type{lua}     & 1         & string,
                                     function  & 在嵌入字形时执行一段 \LUA\ 脚本；如果是函数，则会把
                                                   字体 id 和字符代码传给它 \\
\type{image}   & 1         & image     & 输出一个图像（参数可以是 \type {<image>} 变量，也
                                                   可以是 \type {image_spec} 表） \\
\type{comment} & any       & any       & 此命令的参数被忽略 \\
\bottomrule
\end{longtable}

当字体 id 设为~0 时，它会被替换为当前指定的字体 id。这样就避免了对
未来 id 做取巧处理。通常可以使用 \type {font.nextid}，但如果在此期间
又构建了更复杂的字体，可能已经加载了其他实例。

\type {pdf} 选项还接受一个 \type {mode} 关键字，此时第三个参数设置模式。
该选项会以高效方式更改模式（传入空字符串会在 \PDF\ 文件中多出空行）。
此选项仅对虚拟字体有意义。\type {font} 模式仅在虚拟字体中有意义。各种
模式的定义有些模糊，部分继承自 \PDFTEX。

\begin{longtable}{l p{5cm}}
\toprule 模式           & 说明 \\
\midrule
\type {origin} & 进入页面模式并设置位置 \\
\type {page}   & 进入页面模式 \\
\type {text}   & 进入文本模式 \\
\type {font}   & 进入字体模式（类似文本模式，仅用于虚拟字体） \\
\type {always} & 结束当前字符串，并在需要时强制进行变换 \\
\type {raw}    & 结束当前字符串 \\
\bottomrule
\end{longtable}

你始终需要检查生成的 \PDF\ 代码是什么，因为后端的优化可能产生各种干扰，
而且字体本身就很复杂。下面是一个相当详尽的字形命令示例，使用了这些键：

\begin{typing}
...
commands = {
    { "push" },                     -- remember where we are
    { "right", 5000 },              -- move right about 0.08pt
    { "font", 3 },                  -- select the fonts[3] entry
    { "char", 97 },                 -- place character 97 (ASCII 'a')
 -- { "slot", 2, 97 },              -- an alternative for the previous two
    { "pop" },                      -- go all the way back
    { "down", -200000 },            -- move upwards by about 3pt
    { "special", "pdf: 1 0 0 rg" }  -- switch to red color
 -- { "pdf", "origin", "1 0 0 rg" } -- switch to red color (alternative)
    { "rule", 500000, 20000 }       -- draw a bar
    { "special", "pdf: 0 g" }       -- back to black
 -- { "pdf", "origin", "0 g" }      -- back to black (alternative)
}
...
\end{typing}

在 \type {commands} 数组开始处，\type {font} 的默认值总是~1。因此，
如果虚拟字体本质上只是重编码，通常你不必在数组中显式创建 ``font
命令。

\type {commands} 数组中的矩形只用两个维度构建：它们没有深度。为了正确
的垂直放置，可能需要额外的 \type {down} 命令。

无论你在 \type {commands} 内部产生了多少移动，输出指针始终恰好移动
字符哈希 \type {width} 键所给定的宽度。\type {commands} 数组内部发生的
任何移动在上层都会被忽略。

special 可以带有 \type {pdf:}、\type {pdf:origin:}、\type {pdf:page:}、
\type {pdf:direct:} 或 \type {pdf:raw:} 前缀。当需要拼接字符串时，使用
\type {pdf} 命令可能更高效。

\subsection{人造字体}

即使在 ``真实 字体中，也可以存在虚拟字符。当 \LUATEX\ 在需要
排版某个字符时遇到该字符内部的 \type {commands} 字段，它就会像处理
真正的虚拟字符一样解释这些命令。在这种情况下，如果你没有创建
``fonts 数组，那么默认（也是唯一）的 ``基础 字体就是当前
字体本身。实际上，这意味着你可以为已有字符创建虚拟副本，这在想要创建
复合字符时很有用。

注意：此功能{\it 不\/}能反向使用。虚拟字体中不能有 ``真实 字符！
你也不能用这种技术来做字体重编码；那样做需要一个真正的虚拟字体（因为
已经存在的字符无法被更改）。

\subsection{虚拟字体示例}

\topicindex {fonts+virtual}

最后，下面是一个展示虚拟字体的 plain \TEX\ 输入文件：

\begin{typing}
\directlua {
  callback.register('define_font',
    function (name,size)
      if name == 'cmr10-red' then
        local f = font.read_tfm('cmr10',size)
        f.name  = 'cmr10-red'
        f.type  = 'virtual'
        f.fonts = {
          { name = 'cmr10', size = size }
        }
        for i,v in pairs(f.characters) do
          if string.char(i):find('[tacohanshartmut]') then
             v.commands = {
               { "special", "pdf: 1 0 0 rg" },
               { "char", i },
               { "special", "pdf: 0 g" },
              }
          end
        end
        return f
      else
        return font.read_tfm(name,size)
      end
    end
  )
}

\font\myfont  = cmr10-red at 10pt \myfont  This is a line of text \par
\font\myfontx = cmr10     at 10pt \myfontx Here is another line of text \par
\end{typing}

\section{\type {vf}}

当执行字体 \type {commands} 中定义的 \LUA\ 代码时，可以使用
\type {vf} 库。它提供的函数与命令类似：\type {char}、\type {down}、
\type {fontid}、\type {image}、\type {node}、\type {nop}、\type {pop}、
\type {push}、\type {right}、\nod {rule}、\type {special} 和
\type {pdf}。这个库已经存在一段时间了，但没有怎么宣传和测试，原因之一是
它很容易定义出无效字体（或者搞乱 \PDF\ 流）。请记住，每次输出一个字符时
都会执行这些 \LUA\ 片段。

\section{\type {font}}

\topicindex {fonts+library}

font 库提供了访问字体系统内部结构的接口，还包含加载传统 \TEX\ 字体
度量格式的辅助函数。其他字体加载功能由 \type {fontloader} 库提供，
下一节将讨论它。

\subsection{加载 \TFM\ 文件}

\topicindex {fonts+tfm}

本小节所记录的行为被认为是稳定的，今后不会再有向后不兼容的更改。

\begin{typing}
<table> fnt =
    font.read_tfm(<string> name, <number> s)
\end{typing}

这个数字有点特别：

\begin{itemize}\item 如果为正，它以缩放点指定 ``at size（按指定尺寸）。

\item 如果为负，其绝对值表示相对于字体 designsize 的 ``scaled
    （缩放）设置。

\end{itemize}

\subsection{加载 \VF\ 文件}

\topicindex {fonts+vf}

本小节所记录的行为被认为是稳定的，今后不会再有向后不兼容的更改。

\begin{typing}
<table> vf_fnt =
    font.read_vf(<string> name, <number> s)
\end{typing}

数字 \type {s} 的含义和返回表的格式与 \type {read_tfm} 函数中的类似。

\subsection{fonts 数组}

\topicindex {fonts+virtual}

通过一个虚拟数组，\LUA\ 可以访问 \TEX\ 字体的整个表。

\begin{typing}
font.fonts[n] = { ... }
<table> f = font.fonts[n]
\end{typing}

由于这是一个虚拟数组，你不能对它调用 \type {pairs}，但可参见下面的
\type {font.each} 迭代器。

实现这个虚拟数组的两个元表函数是：

\begin{typing}
<table> f = font.getfont(<number> n)
font.setfont(<number> n, <table> f)
\end{typing}

请注意，目前除非你缓存了它，否则每次访问 \type {font.fonts} 或调用
\type {font.getfont} 都会为整个字体创建一个 \LUA\ 表。如果你想要内部
数据的副本，可以使用 \type {font.getcopy}：

\begin{typing}
<table> f = font.getcopy(<number> n)
\end{typing}

这个函数返回一个 \TEX\ 所知参数的表。这些参数可能与缓存表中的不同：

\begin{typing}
<table> p = font.getparameters(<number> n)
\end{typing}

还要注意：只能对那些已在 \TEX\ 中定义、但自定义以来{\it 完全\/}未被
访问过的字体进行赋值。这大大限制了写访问 \type {font.fonts} 的可用性，
今后很可能会实现一套较宽松的规则。

\subsection{检查字体状态}

你可以通过调用此函数来测试字体的状态：

\begin{typing}
<boolean> f =
    font.frozen(<number> n)
\end{typing}

返回值为 \type {true}（不可赋值）、\type {false}（可以更改）或
\type {nil}（根本不是有效字体）之一。

\subsection{直接定义字体}

\topicindex {fonts+define}

你可以通过调用此函数把自己的字体定义到 \type {font.fonts} 中：

\begin{typing}
<number> i =
    font.define(<table> f)
\end{typing}

返回值是所定义字体的内部 id 编号（即 \type {font.fonts} 中的索引）。
如果创建字体失败，会抛出错误。该表是一个字体结构。另一种调用形式是：

\begin{typing}
<number> i =
    font.define(<number> n, <table> f)
\end{typing}

其中第一个参数是预留的字体 id（见下文）。

\subsection{Unicode 映射}

\type {glyphtounicode} 扩展与 \PDFTEX\ 中的类似，但不提供按字体设置的
方法。这与我们存储字体的方式不同、以及可以重载默认字体加载器有关。如果你
想调整某个特定字体，可以像下面这样在 plain \TEX\ 代码中使用一个
\LUA\ 辅助函数。

\begin{typing}
% ToUnicode vectors will be added

\pdfvariable gentounicode 1

% Wrap a helper in a macro:

\def\luaglyphtounicode#1#2%
  {\directlua{font.settounicode(\fontid\font,\number#1,"#2")}}

% Use the global mapper:

\pdfextension glyphtounicode {A}{0042}

% Test this:

{A \font\test=cmss10 \test  A}

% Define a non unicode symbol font:

\font\orn = OrnementsADF at 10pt

% Register mappings:

\{ \orn
    \luaglyphtounicode{42}{2B9C}
    \luaglyphtounicode{65}{2B9D}
\}

% Test this:

{\orn \char42 \char65 A} A

% And see if cut-n-paste from the pdf file works.
\end{typing}

当然，把辅助函数包进宏的做法取决于具体的宏包。此功能是为 2026 版添加的，
不应干扰先前的行为。顺带说明：对于 \TYPEONE\ 字体，向 pdf 文件添加向量
只在至少存在一个全局扩展条目时发生，但这种由 \LUA\ 驱动的方式没有这个
限制。这意味着当你使用特定于字体的方法时，它也会为其他 8 位字体触发
向量的生成。

\subsection{扩展字体}

\topicindex {fonts+extend}

在合理范围内，你可以在字体定义之后扩展它。由于有些属性最好保持不变，
扩展仅限于添加字符。

\begin{typing}
font.addcharacters(<number> n, <table> f)
\end{typing}

传入的表可以含有 \type {characters} 字段，它是一个与 define 中所用类似的
（子）表；对于虚拟字体，还可以添加 \type {fonts} 表。\type {characters}
表中定义的字符会被添加（当尚不存在时）或替换已有条目。请记住，替换可能有
副作用，因为某个字符可能已经被使用过。我们不设置限制，而是期望用户自己
小心。（\type {setfont} 辅助函数是一种更彻底的替换方式。）

\subsection{预计的下一个字体 id}

\topicindex {fonts+id}

\begin{typing}
<number> i =
    font.nextid()
\end{typing}

它返回如果在代码流的这个位置执行 \type {font.define} 调用时将会返回的
字体 id 编号。这对于需要引用自身的虚拟字体很有用。如果你传入 \type {true}
作为参数，该 id 会被预留，之后可作为第一个参数传给 \type {font.define}。
在创建复杂虚拟字体时这会很方便。

\begin{typing}
<number> i =
    font.nextid(true)
\end{typing}

\subsection{字体 id}

\topicindex {fonts+id}
\topicindex {fonts+current}

\begin{typing}
<number> i =
    font.id(<string> csname)
\end{typing}

它返回与 \type {csname} 关联的字体 id；如果 \type {csname} 未定义，则
返回 $-1$。

\begin{typing}
<number> i =
    font.max()
\end{typing}

这是 \type {font.fonts} 中已使用的最大索引。

\begin{typing}
<number> i = font.current()
font.current(<number> i)
\end{typing}

它获取或设置当前使用的字体编号。

\subsection{遍历所有字体}

\topicindex {fonts+iterate}

\begin{typing}
for i,v in font.each() do
  ...
end
\end{typing}

这是一个遍历每个已定义 \TEX\ 字体的迭代器。返回的第一个值是
\type {font.fonts} 中的索引，第二个值是字体本身（以 \LUA\ 表形式）。
索引按递增顺序列出，但它们并不总是构成连续数字的数组：在某些情况下序列中
可能有空洞。

\subsection{\type{\glyphdimensionsmode}}

早在 \LUATEX\ 的早期，就决定以反映所施垂直偏移的方式来计算字形的有效
高度和深度。高度加上该偏移，而深度只在偏移大于零时才减去该偏移。现在我们
可以通过这个模式参数更细致地控制这一行为。偏移会加到高度上，和|/|或从
深度中减去。有效值永不为负。零模式是默认值。

\begin{longtable}{l p{8cm}}
\toprule 取值     & 效果 \\
\midrule
\quad\type {0} & 旧行为：把偏移加到高度上，并且仅当偏移为正时才从深度
                  中减去偏移 \\
\type {1} & 把偏移加到高度上，并从深度中减去偏移 \\
\quad\type {2} & 把偏移加到高度上并从深度中减去偏移，但保留当前结果与
                  先前结果的较大值 \\
\type {3} & 使用字形本身的高度和深度，即不应用偏移 \\
\bottomrule
\end{longtable}

\subsection{\type {\discretionaryligaturemode}}

此参数控制复杂连字如何与断字点（由连字符处理程序插入的）交互。非零值会
阻止构造所谓的 init 和 select 断字点。

% ConTeXt \definefontfeature/\definefont 演示已省略。
% 以下为三个模式（0、1、2）的连字与断字交互演示，图已省略。

\subsection{\type{\tracinglostchars}}

这个跟踪器报告缺失的字符。按设计，我们把报告推迟到后端包含它们时、数学
排版需要它们时，或创建虚拟字体时（1、2、3）。应 \LATEX\ 团队的要求，
我们另外增加了两个在更早阶段报告的选项，即在字形被打包或在行断行期间
处理时报告（4 和 5）。

\begin{longtable}{r l}
    \toprule     & 效果 \\
    \midrule
0   & 不报告 \\
1   & 仅向日志文件发出警告 \\
2   & 警告并强制输出到终端 \\
3   & 错误 \\
4   & 提前警告（并强制输出到终端） \\
> 4 & 提前报错 \\
    \bottomrule
\end{longtable}
